% front_matter.tex -- Abstract and Section 1 (Introduction)

\begin{abstract}
As AI agents are deployed in increasingly complex social and economic
settings, ensuring they reason ethically about cooperation, trust, and
fairness becomes critical for alignment. We present \textsc{Kant}, a
comprehensive game-theoretic training and evaluation framework that goes
far beyond simple benchmarking. Built on the OpenEnv platform, Kant
provides a library of 114 games spanning ten strategic domains---from
classical dilemmas through information games, market competition, and
auctions to cooperative game theory and adaptive games with
history-dependent payoffs. The framework supports $N$-player
interactions, coalition formation with negotiation and enforcement
mechanisms, a novel meta-governance engine that allows agents to
\emph{change the rules of the game during play} through democratic
proposals and voting, and a reputation system with cross-episode social
memory backed by knowledge graph queries. Dynamic game creation enables
open-ended evaluation
where agents construct new strategic environments at runtime. We define a
training pipeline using GRPO, DPO, and self-play on game trajectories
with curriculum learning across game domains, and evaluate safety
transfer on external
benchmarks including HarmBench, ETHICS, TruthfulQA, XSTest, and MT-Bench.
By combining comprehensive game-theoretic environments with alignment
training, Kant provides a platform for studying whether strategic reasoning
can serve as a foundation for teaching ethical behavior to language models.
\end{abstract}

\section{Introduction}
\label{sec:intro}

The deployment of large language model (LLM) agents in real-world
settings---negotiation, resource allocation, market participation, collective
governance---raises a pressing question: \emph{can we systematically train
these agents to reason ethically about strategic interactions?} Classical game
theory offers a rich library of settings that distil key social dilemmas into
tractable formal
games~\citep{nash1950equilibrium,rapoport1965prisoner,axelrod1984evolution}.
Yet existing AI benchmarks either focus on single-agent decision-making
or embed strategic reasoning inside complex, high-dimensional environments
that make controlled evaluation
difficult~\citep{pan2023machiavelli,leibo2021meltingpot}.

We present \textsc{Kant}, a framework that uses game theory as an
\emph{alignment substrate}---a structured environment in which ethical
reasoning can be trained, evaluated, and understood. Rather than treating
game-theoretic evaluation as an end in itself, Kant asks whether the
strategic competencies developed through game play---cooperation, fairness,
trust calibration, exploitation resistance---transfer to broader alignment
objectives.

The framework makes several contributions:

\begin{enumerate}[nosep]
    \item A library of 114 games spanning ten strategic domains, from
          classical dilemmas to auctions, cooperative games, adaptive
          games, and contests (Section~\ref{sec:game-library}).
    \item $N$-player extensions and a coalition formation system with
          negotiation, enforcement, and side payments
          (Sections~\ref{sec:nplayer}--\ref{sec:coalition}).
    \item A meta-governance engine enabling democratic rule modification
          during play, connecting game theory to Constitutional AI
          (Section~\ref{sec:governance}).
    \item A reputation and social memory system enabling cross-episode
          learning via gossip variants and knowledge graph queries
          (Section~\ref{sec:reputation}).
    \item Dynamic game creation for open-ended evaluation
          (Section~\ref{sec:dynamic}).
    \item Five adaptive game factories with history-dependent payoffs
          that evolve during play (Section~\ref{sec:adaptive}).
    \item A training pipeline using GRPO, DPO, and self-play on game
          trajectories with curriculum learning
          (Section~\ref{sec:training}).
    \item Safety transfer evaluation on external benchmarks
          (Section~\ref{sec:safety}).
    \item An alignment-oriented evaluation protocol with six metrics
          and stratified game splitting
          (Section~\ref{sec:eval-protocol}).
\end{enumerate}
